
\section{Conclusion and Comparison of PRNG Methods}

This section summarizes and compares the pseudorandom number generators discussed in the previous chapters. The goal is to highlight their main differences in terms of period length, computational efficiency, statistical quality, and structural complexity.

\subsection{General Comparison Criteria}

Pseudorandom number generators can be evaluated according to several fundamental criteria:

\begin{itemize}
    \item \textbf{Period length:} the length of the cycle before the sequence repeats.
    \item \textbf{Computational efficiency:} the cost of generating a single value.
    \item \textbf{Statistical quality:} how well the output approximates a truly random sequence.
    \item \textbf{Structural complexity:} mathematical and implementation complexity of the generator.
\end{itemize}

These criteria are often in conflict; for example, higher statistical quality typically comes at the cost of increased complexity. This is why practical evaluation usually combines mathematical period analysis, implementation-level considerations and empirical tests such as histogram checks, bit-level statistics and more advanced test suites~\cite{knuth1997art,lecuyer2007testu01}.

\subsection{Classical Generators}

Classical methods such as linear congruential generators (LCG) and Lehmer generators are based on simple modular arithmetic:

\[
x_{n+1} = (ax_n + c) \bmod m
\]

or in the multiplicative case:

\[
x_{n+1} = ax_n \bmod m.
\]

These generators are computationally very efficient and easy to implement. However, their statistical properties depend strongly on parameter choice. The period of an LCG is at most $m$, and full period is achieved only under the Hull--Dobell conditions~\cite{hull1962random}. A known drawback is their tendency to exhibit geometric structure when viewed in higher dimensions, which reduces their suitability for high-quality simulations~\cite{lecuyer1999tables}.

\subsection{Recursive and Shift-Based Generators}

More advanced recursive methods such as lagged Fibonacci generators (LFG) and linear feedback shift registers (LFSR) improve statistical properties by introducing longer memory into the state:

\begin{itemize}
    \item LFG uses previous values in a recurrence relation.
    \item LFSR operates on bit-level linear recurrences over finite fields.
\end{itemize}

These methods generally achieve longer periods and better distribution properties than simple LCGs, while maintaining moderate computational cost. However, as shown in the experiments, poor taps or weak initialization can still create short cycles or visible structure. This makes parameter selection essential, especially for linear generators over finite fields~\cite{golomb1967shift}.

\subsection{Modern High-Quality Generators}

Modern generators such as the Mersenne Twister, PCG, and xorshift-based methods are designed to optimize statistical quality and period length.

The Mersenne Twister achieves an extremely large period of $2^{19937}-1$, making repetition practically impossible in most ordinary applications~\cite{matsumoto1998mersenne}. PCG improves statistical behavior by applying permutation functions to linear congruential states, while maintaining efficiency~\cite{oneill2014pcg}. Xorshift generators rely on bitwise XOR and shift operations, providing extremely fast generation with good statistical performance when properly parameterized~\cite{marsaglia2003xorshift}.

\subsection{Trade-offs Between Methods}

There is no single optimal PRNG for all applications. Instead, different methods optimize different aspects:

\begin{itemize}
    \item \textbf{LCG / Lehmer:} very fast, but weaker statistical quality and strong dependence on parameters.
    \item \textbf{LFSR / LFG:} better structure and longer memory, but sensitive to taps and initialization.
    \item \textbf{Mersenne Twister:} excellent general-purpose statistical properties and very long period, but larger state complexity.
    \item \textbf{PCG / Xorshift:} modern balance between speed, small implementation size and practical statistical quality.
\end{itemize}

\subsection{Final Remarks}

Pseudorandom number generation is fundamentally a trade-off between simplicity, speed, and statistical quality. While classical methods such as LCG provide insight into the structure of deterministic sequences, modern generators achieve practical randomness suitable for large-scale simulations and computational applications.

The experiments in this work also show that visual plots and simple statistics are useful, but not sufficient on their own. A generator may look acceptable in a histogram while still having a poor period, bad bit-level behavior or structural correlations. Therefore, understanding the underlying mathematical structure is essential for selecting appropriate methods in scientific computing, cryptography-related contexts and numerical simulations.
